\chapter{Channels and \texttt{select}}\label{ch:channels-select}

\index{Channel@\texttt{Channel}}
\index{channels}

In the previous chapter, we used channels as convenient mailboxes---toss a value in one end, pull it out the other.
But channels in Lattice are a full concurrency primitive: they are the primary mechanism for safe communication between threads, and with \lstinline|select|, they become a powerful multiplexing tool that lets you wait on multiple sources of data simultaneously.

Let's start from the beginning.

%% ─────────────────────────────────────────────────────────────
\section{\texttt{Channel::new()}, \texttt{send()}, \texttt{recv()}, \texttt{close()}}\label{sec:channel-basics}
%% ─────────────────────────────────────────────────────────────

\index{Channel@\texttt{Channel}!creation}
\index{Channel::new@\texttt{Channel::new()}}

A channel is a thread-safe FIFO queue.
You create one with \lstinline|Channel::new()|, send values into it, and receive values from it:

\begin{lstlisting}[language=Lattice, caption={Channel basics}]
let ch = Channel::new()

ch.send(42)
ch.send("hello")
ch.send(true)

print(ch.recv())  // 42
print(ch.recv())  // hello
print(ch.recv())  // true
\end{lstlisting}

Channels are untyped---you can send any Lattice value through them.
Values come out in the order they were sent (first-in, first-out).

\begin{defnbox}{Channel}
\index{Channel@\texttt{Channel}!definition}
A \emph{Channel} is an unbounded, thread-safe FIFO queue.
It supports three core operations: \lstinline|send(value)| enqueues a value, \lstinline|recv()| dequeues a value (blocking if the channel is empty), and \lstinline|close()| signals that no more values will be sent.
Internally, a channel is a mutex-protected vector with a condition variable for blocking receives (see \filename{include/channel.h}).
\end{defnbox}

\subsection{\texttt{send()} --- Putting Values In}

\index{Channel@\texttt{Channel}!send}
\index{send@\texttt{send()}}

The \lstinline|send()| method enqueues a value into the channel's buffer:

\begin{lstlisting}[language=Lattice, caption={Sending values}]
let orders = Channel::new()

orders.send("latte")
orders.send("cappuccino")
orders.send("espresso")

print("Queued ${3} orders")
\end{lstlisting}

Sending is non-blocking: because Lattice channels are unbounded, \lstinline|send()| always succeeds immediately (as long as the channel is open).
Under the hood, \texttt{channel\_send} in \filename{src/channel.c} acquires the channel's mutex, pushes the value onto the internal \texttt{LatVec} buffer, signals the \texttt{cond\_notempty} condition variable to wake any blocked receivers, and also wakes any select waiters.

\begin{notebox}{Sending to a Closed Channel}
\index{Channel@\texttt{Channel}!closed}
If you \lstinline|send()| to a closed channel, the value is silently dropped and the send returns \lstinline|false| internally.
The value is freed to prevent memory leaks.
Design your programs so that producers stop sending before or shortly after the channel is closed.
\end{notebox}

\subsection{\texttt{recv()} --- Taking Values Out}

\index{Channel@\texttt{Channel}!recv}
\index{recv@\texttt{recv()}}

The \lstinline|recv()| method dequeues the next value from the channel:

\begin{lstlisting}[language=Lattice, caption={Receiving values}]
let ch = Channel::new()
ch.send(100)
ch.send(200)

let first = ch.recv()   // 100
let second = ch.recv()  // 200
print("Got ${first} and ${second}")
\end{lstlisting}

Here is the critical behavior: \textbf{if the channel is empty and still open, \lstinline|recv()| blocks the calling thread until a value is available}.
This blocking behavior is what makes channels useful for synchronization---a consumer naturally waits for its producer.

\begin{lstlisting}[language=Lattice, caption={Blocking recv}]
let ch = Channel::new()

scope {
  spawn {
    // Simulate slow work
    let start = clock()
    while clock() - start < 100 { }
    ch.send("result ready")
  }
  spawn {
    // This recv blocks until the value arrives
    let result = ch.recv()
    print("Got: ${result}")
  }
}
\end{lstlisting}

In the implementation (\filename{src/channel.c}, \texttt{channel\_recv}), receiving works by locking the mutex and entering a \texttt{while} loop that calls \texttt{pthread\_cond\_wait} as long as the buffer is empty and the channel is not closed.
When a value is available, it is shifted from the front of the buffer (FIFO semantics) and returned.

\subsection{What Happens When You \texttt{recv()} from a Closed, Empty Channel?}

If the channel is closed and its buffer is empty, \lstinline|recv()| returns \lstinline|nil|.
This provides a clean way to detect that a producer is done:

\begin{lstlisting}[language=Lattice, caption={Receiving from a closed channel}]
let ch = Channel::new()
ch.send(1)
ch.send(2)
ch.close()

print(ch.recv())  // 1
print(ch.recv())  // 2
print(ch.recv())  // nil (channel closed and empty)
\end{lstlisting}

\subsection{\texttt{close()} --- Signaling Completion}

\index{Channel@\texttt{Channel}!close}
\index{close@\texttt{close()}}

The \lstinline|close()| method signals that no more values will be sent:

\begin{lstlisting}[language=Lattice, caption={Closing a channel}]
let tasks = Channel::new()

tasks.send("task_a")
tasks.send("task_b")
tasks.close()

// Receivers can still drain buffered values
print(tasks.recv())  // task_a
print(tasks.recv())  // task_b
print(tasks.recv())  // nil
\end{lstlisting}

Closing is a one-way operation---once closed, a channel stays closed.
The implementation (\texttt{channel\_close} in \filename{src/channel.c}) sets the \texttt{closed} flag, then broadcasts on the condition variable to wake all blocked receivers and select waiters.
This ensures no thread stays blocked forever on a closed channel.

\begin{tipbox}{Always Close Producer Channels}
\index{Channel@\texttt{Channel}!closing best practice}
When a producer is done sending, close the channel.
This lets consumers detect the end of the stream and exit their loops cleanly.
A channel that is never closed will cause any \lstinline|recv()| call to block indefinitely once the buffer is drained.
\end{tipbox}

\subsection{Channel Lifecycle}

Channels are reference-counted.
When you share a channel across spawns, each copy increments the reference count (\texttt{channel\_retain}).
When a copy goes out of scope, the count is decremented (\texttt{channel\_release}).
The channel's internal buffer is freed only when the last reference is released.

\begin{lstlisting}[language=Lattice, caption={Channel lifecycle across spawns}]
fn producer(ch: Channel, items: Array) {
  for item in items {
    ch.send(item)
  }
  ch.close()
}

fn consumer(ch: Channel) {
  flux item = ch.recv()
  while item != nil {
    print("Consumed: ${item}")
    item = ch.recv()
  }
}

let pipeline = Channel::new()

scope {
  spawn { producer(pipeline, ["apple", "banana", "cherry"]) }
  spawn { consumer(pipeline) }
}
// Output:
// Consumed: apple
// Consumed: banana
// Consumed: cherry
\end{lstlisting}

%% ─────────────────────────────────────────────────────────────
\section{Producer-Consumer Patterns}\label{sec:producer-consumer}
%% ─────────────────────────────────────────────────────────────

\index{producer-consumer}
\index{concurrency patterns!producer-consumer}

The producer-consumer pattern is the bread and butter of channel-based concurrency.
One thread produces data, another consumes it, and the channel decouples them.

\subsection{Single Producer, Single Consumer}

The most common variant:

\begin{lstlisting}[language=Lattice, caption={Single producer, single consumer}]
let work_queue = Channel::new()

scope {
  // Producer
  spawn {
    for i in 0..10 {
      work_queue.send(i * i)
    }
    work_queue.close()
  }

  // Consumer
  spawn {
    flux val = work_queue.recv()
    while val != nil {
      print("Processing: ${val}")
      val = work_queue.recv()
    }
    print("Consumer done")
  }
}
\end{lstlisting}

The consumer loop is idiomatic: call \lstinline|recv()|, check for \lstinline|nil| (which signals the channel is closed and drained), and process each value.

\subsection{Multiple Producers, Single Consumer}

Multiple spawns can send to the same channel:

\begin{lstlisting}[language=Lattice, caption={Multiple producers, single consumer}]
let log_channel = Channel::new()

scope {
  spawn {
    for i in 0..5 {
      log_channel.send("[Server A] Request ${i}")
    }
  }
  spawn {
    for i in 0..5 {
      log_channel.send("[Server B] Request ${i}")
    }
  }
  spawn {
    // Give producers a moment, then close
    let start = clock()
    while clock() - start < 50 { }
    log_channel.close()
  }
}

// Drain the log
flux msg = log_channel.recv()
while msg != nil {
  print(msg)
  msg = log_channel.recv()
}
\end{lstlisting}

Because the channel is mutex-protected, sends from different threads are serialized automatically---no data corruption, no lost messages.

\subsection{Single Producer, Multiple Consumers}

This pattern distributes work across consumers:

\begin{lstlisting}[language=Lattice, caption={Work distribution across consumers}]
fn worker(id: Int, jobs: Channel, results: Channel) {
  flux job = jobs.recv()
  while job != nil {
    let result = "Worker ${id} processed job ${job}"
    results.send(result)
    job = jobs.recv()
  }
}

let jobs = Channel::new()
let results = Channel::new()

// Enqueue work
for i in 0..12 {
  jobs.send(i)
}
jobs.close()

scope {
  spawn { worker(1, jobs, results) }
  spawn { worker(2, jobs, results) }
  spawn { worker(3, jobs, results) }
}

// Collect all results
for _ in 0..12 {
  print(results.recv())
}
\end{lstlisting}

Each worker calls \lstinline|recv()| on the same jobs channel.
Because \lstinline|recv()| is atomic (only one thread gets each value), the work is naturally distributed---no explicit load balancing needed.

%% ─────────────────────────────────────────────────────────────
\section{\texttt{select} with Multiple Channels}\label{sec:select-channels}
%% ─────────────────────────────────────────────────────────────

\index{select@\texttt{select}}
\index{channel multiplexing}

What if you need to wait on \emph{multiple} channels at once?
Maybe you have two data sources and want to process whichever one has data ready first.
This is where \lstinline|select| comes in.

\begin{lstlisting}[language=Lattice, caption={Basic select}]
let weather = Channel::new()
let news = Channel::new()

scope {
  spawn { weather.send("Sunny, 72F") }
  spawn { news.send("Lattice 1.0 released!") }
}

let result = select {
  weather -> report { "Weather: ${report}" }
  news -> headline { "News: ${headline}" }
}
print(result)
// Either "Weather: Sunny, 72F" or "News: Lattice 1.0 released!"
\end{lstlisting}

A \lstinline|select| expression waits on multiple channels and executes the arm corresponding to the first channel that has data available.

\begin{defnbox}{\texttt{select} Expression}
\index{select@\texttt{select}!definition}
A \lstinline|select| expression monitors multiple channels simultaneously and executes the body of the first arm that successfully receives a value.
If multiple channels have data ready, one is chosen at random (via Fisher-Yates shuffle in the VM) to ensure fairness.
The expression evaluates to the result of the chosen arm's body.
\end{defnbox}

\subsection{Anatomy of a \texttt{select} Arm}

Each arm of a \lstinline|select| has three parts:

\begin{lstlisting}[language=Lattice, caption={Select arm anatomy}]
select {
  channel_expr -> binding_name { body }
}
\end{lstlisting}

\begin{itemize}
  \item \texttt{channel\_expr}: An expression that evaluates to a \lstinline|Channel|.
  \item \texttt{binding\_name}: A variable name that will hold the received value inside the body. This is optional---you can omit the binding if you don't need the value.
  \item \texttt{body}: The code to execute when this channel delivers a value.
\end{itemize}

\subsection{Fairness}

\index{select@\texttt{select}!fairness}

When multiple channels have data ready simultaneously, Lattice does not always pick the first arm in source order.
Instead, the VM shuffles the arm indices using the Fisher-Yates algorithm before polling.
This prevents starvation: a fast producer on one channel cannot permanently starve a slower producer on another.

\begin{lstlisting}[language=Lattice, caption={Fairness in select}]
let fast = Channel::new()
let slow = Channel::new()

// Load both channels
for i in 0..10 {
  fast.send("fast-${i}")
  slow.send("slow-${i}")
}

// Each select picks fairly from both
for _ in 0..5 {
  let result = select {
    fast -> msg { msg }
    slow -> msg { msg }
  }
  print(result)
}
// You'll see a mix of "fast-*" and "slow-*" messages
\end{lstlisting}

\subsection{Select in a Loop}

A common pattern is to use \lstinline|select| inside a loop to continuously process events from multiple sources:

\begin{lstlisting}[language=Lattice, caption={Event loop with select}]
let keyboard = Channel::new()
let network = Channel::new()
let timer = Channel::new()

// Simulate event sources
scope {
  spawn {
    keyboard.send("key_press: Enter")
    keyboard.send("key_press: Escape")
    keyboard.close()
  }
  spawn {
    network.send("packet: 200 OK")
    network.close()
  }
  spawn {
    timer.send("tick")
    timer.send("tick")
    timer.close()
  }
}

// Process all events
flux running = true
while running {
  let event = select {
    keyboard -> key { "Keyboard: ${key}" }
    network -> pkt { "Network: ${pkt}" }
    timer -> t { "Timer: ${t}" }
    default { running = false; "done" }
  }
  print(event)
}
\end{lstlisting}

The \lstinline|default| arm fires when all channels are empty (or closed), allowing us to break out of the loop cleanly.

%% ─────────────────────────────────────────────────────────────
\section{\texttt{default} and \texttt{timeout} Arms}\label{sec:select-default-timeout}
%% ─────────────────────────────────────────────────────────────

\index{select@\texttt{select}!default arm}
\index{select@\texttt{select}!timeout arm}

Plain \lstinline|select| blocks until one of its channel arms can receive.
Two special arms modify this behavior.

\subsection{The \texttt{default} Arm}

The \lstinline|default| arm executes immediately if no channel has data ready:

\begin{lstlisting}[language=Lattice, caption={Select with default}]
let ch = Channel::new()

let result = select {
  ch -> val { "Got: ${val}" }
  default { "Nothing available" }
}
print(result)  // Nothing available
\end{lstlisting}

This turns \lstinline|select| into a non-blocking check.
The \lstinline|default| arm is useful for polling patterns where you want to do other work when no messages are ready.

\begin{lstlisting}[language=Lattice, caption={Polling pattern with default}]
let inbox = Channel::new()

scope {
  spawn {
    let start = clock()
    while clock() - start < 100 { }
    inbox.send("delayed message")
  }
  spawn {
    flux attempts = 0
    flux found = false
    while found == false {
      let result = select {
        inbox -> msg {
          found = true
          msg
        }
        default {
          attempts = attempts + 1
          "waiting..."
        }
      }
      if found {
        print("Got message after ${attempts} polls: ${result}")
      }
    }
  }
}
\end{lstlisting}

\begin{warnbox}{Busy-Waiting with Default}
A \lstinline|select| with a \lstinline|default| arm never blocks.
If you put it in a tight loop, you will burn CPU cycles polling.
Consider using a \lstinline|timeout| arm instead if you can tolerate some latency.
\end{warnbox}

\subsection{The \texttt{timeout} Arm}

\index{timeout@\texttt{timeout}!in select}

The \lstinline|timeout| arm takes a duration in milliseconds and fires if no channel delivers a value within that time:

\begin{lstlisting}[language=Lattice, caption={Select with timeout}]
let slow_service = Channel::new()

let result = select {
  slow_service -> data { "Got: ${data}" }
  timeout 500 { "Timed out after 500ms" }
}
print(result)  // Timed out after 500ms
\end{lstlisting}

The timeout is implemented using \texttt{pthread\_cond\_timedwait} in the VM (see \filename{src/stackvm.c}, the \texttt{OP\_SELECT} handler).
The VM computes an absolute deadline from the current time plus the timeout value, then waits on the shared condition variable with that deadline.

\begin{lstlisting}[language=Lattice, caption={Timeout with fallback logic}]
fn fetch_with_timeout(source: Channel, timeout_ms: Int) {
  let result = select {
    source -> data { data }
    timeout timeout_ms { nil }
  }
  if result == nil {
    print("Request timed out, using cached data")
    return "cached_value"
  }
  return result
}

let api = Channel::new()

// Simulate a slow API
scope {
  spawn {
    let start = clock()
    while clock() - start < 1000 { }
    api.send("fresh data")
  }
  spawn {
    let value = fetch_with_timeout(api, 200)
    print("Using: ${value}")
  }
}
// Output: Request timed out, using cached data
//         Using: cached_value
\end{lstlisting}

\subsection{Combining \texttt{default} and \texttt{timeout}}

You can have \emph{either} a \lstinline|default| or a \lstinline|timeout| arm, but not both.
If a \lstinline|default| arm is present, it always takes priority over any timeout since it fires immediately when no data is ready.

\begin{tipbox}{Choosing Between \texttt{default} and \texttt{timeout}}
\begin{itemize}
  \item Use \lstinline|default| when you want to do other work \emph{immediately} if no message is ready (non-blocking poll).
  \item Use \lstinline|timeout| when you can afford to wait a bit but need an upper bound on how long (bounded blocking).
  \item Use neither when you are happy to wait indefinitely for a message (unbounded blocking).
\end{itemize}
\end{tipbox}

\subsection{Select with All Channels Closed}

If all channel arms refer to closed, empty channels:

\begin{itemize}
  \item If a \lstinline|default| arm is present, it executes.
  \item If a \lstinline|timeout| arm is present, the select waits until the timeout, then executes the timeout arm.
  \item If neither is present, the select unblocks and returns \lstinline|Unit|.
\end{itemize}

\begin{lstlisting}[language=Lattice, caption={Select with closed channels}]
let ch = Channel::new()
ch.close()

let result = select {
  ch -> val { "Got: ${val}" }
  default { "All channels closed" }
}
print(result)  // All channels closed
\end{lstlisting}

%% ─────────────────────────────────────────────────────────────
\section{Building Pipelines with Channels}\label{sec:channel-pipelines}
%% ─────────────────────────────────────────────────────────────

\index{pipeline pattern}
\index{concurrency patterns!pipeline}

One of the most elegant uses of channels is building data processing pipelines, where each stage reads from one channel, transforms the data, and writes to another.

\subsection{A Two-Stage Pipeline}

\begin{lstlisting}[language=Lattice, caption={Two-stage channel pipeline}]
fn generate(output: Channel) {
  for i in 1..11 {
    output.send(i)
  }
  output.close()
}

fn square(input: Channel, output: Channel) {
  flux val = input.recv()
  while val != nil {
    output.send(val * val)
    val = input.recv()
  }
  output.close()
}

let stage1 = Channel::new()
let stage2 = Channel::new()

scope {
  spawn { generate(stage1) }
  spawn { square(stage1, stage2) }
}

// Consume the final output
flux result = stage2.recv()
while result != nil {
  print(result)
  result = stage2.recv()
}
// Output: 1, 4, 9, 16, 25, 36, 49, 64, 81, 100
\end{lstlisting}

Each stage runs concurrently.
As soon as the generator produces a value, the squarer can consume and transform it---the pipeline processes data in a streaming fashion rather than waiting for all data to be generated first.

\subsection{A Three-Stage Pipeline}

Let's add a filtering stage:

\begin{lstlisting}[language=Lattice, caption={Three-stage pipeline: generate, filter, format}]
fn generate_numbers(output: Channel, limit: Int) {
  for i in 1..limit {
    output.send(i)
  }
  output.close()
}

fn filter_evens(input: Channel, output: Channel) {
  flux val = input.recv()
  while val != nil {
    if val % 2 == 0 {
      output.send(val)
    }
    val = input.recv()
  }
  output.close()
}

fn format_output(input: Channel, output: Channel) {
  flux val = input.recv()
  while val != nil {
    output.send("Even number: ${val}")
    val = input.recv()
  }
  output.close()
}

let raw = Channel::new()
let filtered = Channel::new()
let formatted = Channel::new()

scope {
  spawn { generate_numbers(raw, 21) }
  spawn { filter_evens(raw, filtered) }
  spawn { format_output(filtered, formatted) }
}

flux msg = formatted.recv()
while msg != nil {
  print(msg)
  msg = formatted.recv()
}
// Output:
// Even number: 2
// Even number: 4
// ...
// Even number: 20
\end{lstlisting}

\subsection{Fan-Out Pipeline}

You can split a pipeline by having multiple consumers read from the same channel:

\begin{lstlisting}[language=Lattice, caption={Fan-out: two workers consuming one channel}]
fn heavy_compute(id: Int, input: Channel, output: Channel) {
  flux val = input.recv()
  while val != nil {
    // Simulate computation
    output.send("Worker ${id}: ${val * val}")
    val = input.recv()
  }
}

let input = Channel::new()
let output = Channel::new()

for i in 1..21 {
  input.send(i)
}
input.close()

scope {
  spawn { heavy_compute(1, input, output) }
  spawn { heavy_compute(2, input, output) }
}

for _ in 0..20 {
  print(output.recv())
}
\end{lstlisting}

Work is distributed across the two workers automatically because each \lstinline|recv()| is atomic---only one worker gets each value.

%% ─────────────────────────────────────────────────────────────
\section{Async Iterators: \texttt{async\_iter}, \texttt{async\_map}, \texttt{async\_filter}}\label{sec:async-iterators}
%% ─────────────────────────────────────────────────────────────

\index{async\_iter@\texttt{async\_iter}}
\index{async\_map@\texttt{async\_map}}
\index{async\_filter@\texttt{async\_filter}}
\index{iterators!async}

Building pipelines with raw channels is powerful but verbose.
Lattice provides three built-in functions that bridge channels and iterators, giving you a more ergonomic way to build concurrent data flows.

\subsection{\texttt{async\_iter} --- Channel-Backed Iterators}

\lstinline|async_iter| takes a closure that receives a channel, spawns the closure on a background thread, and returns an iterator that yields values from that channel:

\begin{lstlisting}[language=Lattice, caption={Creating an async iterator}]
let numbers = async_iter(|ch| {
  for i in 1..6 {
    ch.send(i * 10)
  }
  // Channel is automatically closed when the closure returns
})

for val in numbers {
  print(val)
}
// Output:
// 10
// 20
// 30
// 40
// 50
\end{lstlisting}

\begin{defnbox}{\texttt{async\_iter}}
\index{async\_iter@\texttt{async\_iter}!definition}
\lstinline|async_iter(producer_fn)| creates a channel, spawns a detached thread running \lstinline|producer_fn| with the channel as its argument, and returns an iterator that pulls values from the channel.
When the producer closure returns, the channel is automatically closed, which causes the iterator to terminate.
\end{defnbox}

Under the hood (\filename{src/runtime.c}, \texttt{native\_async\_iter}), the function creates a \texttt{LatChannel}, clones the current VM for thread safety, and launches a detached thread that calls the closure with the channel.
The returned iterator wraps the channel and calls \texttt{channel\_recv} on each \texttt{next()}.

\subsection{\texttt{async\_map} --- Transform an Async Stream}

\lstinline|async_map| applies a transformation function to each element of an iterator:

\begin{lstlisting}[language=Lattice, caption={Async map transformation}]
let raw_data = async_iter(|ch| {
  for i in 1..6 {
    ch.send(i)
  }
})

let doubled = async_map(raw_data, |x| x * 2)

for val in doubled {
  print(val)
}
// Output: 2, 4, 6, 8, 10
\end{lstlisting}

The transformation runs synchronously on the consumer thread as values are pulled---this is \emph{not} an additional concurrent stage.
If you need the transformation itself to run concurrently, use an explicit channel pipeline as shown in \cref{sec:channel-pipelines}.

\subsection{\texttt{async\_filter} --- Filter an Async Stream}

\lstinline|async_filter| keeps only elements that pass a predicate:

\begin{lstlisting}[language=Lattice, caption={Async filter}]
let source = async_iter(|ch| {
  for i in 1..21 {
    ch.send(i)
  }
})

let evens = async_filter(source, |x| x % 2 == 0)

for val in evens {
  print(val)
}
// Output: 2, 4, 6, 8, 10, 12, 14, 16, 18, 20
\end{lstlisting}

\subsection{Composing Async Operations}

These functions compose naturally into a pipeline:

\begin{lstlisting}[language=Lattice, caption={Composed async pipeline}]
let pipeline = async_iter(|ch| {
  for i in 1..101 {
    ch.send(i)
  }
})

let result = async_filter(
  async_map(pipeline, |x| x * x),
  |x| x % 3 == 0
)

flux count = 0
for val in result {
  count = count + 1
}
print("Found ${count} perfect squares divisible by 3")
\end{lstlisting}

This creates a three-stage processing flow: generate numbers 1--100, square them, then keep only those divisible by 3.
The generation runs on a background thread; the mapping and filtering run on the consumer thread as values are pulled through.

\begin{tipbox}{When to Use Async Iterators vs.\ Raw Channels}
\index{async iterators!vs channels}
Use \lstinline|async_iter| when you have a single producer and want to consume its output with familiar \lstinline|for| loop syntax.
Use raw channels when you need multiple producers, multiple consumers, bidirectional communication, or explicit \lstinline|select| multiplexing.
\end{tipbox}

%% ─────────────────────────────────────────────────────────────
\section{Under the Hood: \texttt{select} Implementation}\label{sec:select-internals}
%% ─────────────────────────────────────────────────────────────

\index{select@\texttt{select}!implementation}
\index{OP\_SELECT@\texttt{OP\_SELECT}}

For those curious about the machinery, here is how \lstinline|select| works inside the VM.

\subsection{Compilation}

The compiler (\filename{src/stackcompiler.c}, \texttt{EXPR\_SELECT}) emits a single \lstinline|OP_SELECT| instruction followed by per-arm metadata:

\begin{enumerate}
  \item \textbf{Arm count} (1 byte): how many arms the select has.
  \item For each arm, four bytes:
    \begin{itemize}
      \item \textbf{Flags}: bit 0 = default, bit 1 = timeout, bit 2 = has binding name.
      \item \textbf{Channel/timeout index}: index of a sub-chunk constant that evaluates to the channel (or timeout value).
      \item \textbf{Body index}: index of the sub-chunk for the arm's body.
      \item \textbf{Binding index}: index of the string constant for the binding name (or \texttt{0xFF} if no binding).
    \end{itemize}
\end{enumerate}

\subsection{Execution}

The VM's \lstinline|OP_SELECT| handler (\filename{src/stackvm.c}) follows this algorithm:

\begin{enumerate}
  \item Export locals to a new environment scope (so sub-chunks can see them).
  \item Evaluate all channel expressions to obtain \texttt{LatChannel*} pointers.
  \item Evaluate the timeout expression (if present) to get a millisecond value.
  \item Build a shuffled index array (Fisher-Yates) for fairness.
  \item Enter a loop:
    \begin{enumerate}
      \item Try a non-blocking \texttt{channel\_try\_recv} on each channel arm (in shuffled order).
      \item If a value is received: bind it in the environment, execute the arm body, and break.
      \item If all channels are closed: execute the default arm (if present) and break.
      \item If a default arm exists: execute it immediately (non-blocking) and break.
      \item Otherwise: register a \texttt{LatSelectWaiter} on all channels, then \texttt{pthread\_cond\_wait} (or \texttt{pthread\_cond\_timedwait} if there is a timeout).
      \item On wake: remove waiters and retry from step (a).
    \end{enumerate}
  \item Clean up: release channels, pop scope, push result.
\end{enumerate}

The \texttt{LatSelectWaiter} structure (defined in \filename{include/channel.h}) is a linked-list node containing a shared mutex and condition variable.
When any monitored channel receives data or is closed, it signals the waiter's condition variable, waking the select loop.

%% ─────────────────────────────────────────────────────────────
\section{Exercises}\label{sec:ch20-exercises}
%% ─────────────────────────────────────────────────────────────

\begin{enumerate}[label=\textbf{\arabic*.}]

  \item \textbf{Ping-Pong.} Create two channels and two spawns that play ping-pong: spawn A sends ``ping'' on channel 1 and waits for ``pong'' on channel 2; spawn B does the reverse. Run 10 rounds and print each exchange.

  \item \textbf{Timeout Handler.} Write a function that takes a channel and a timeout in milliseconds. Using \lstinline|select|, return the first value from the channel or the string \lstinline|"timeout"| if the deadline passes. Test with both fast and slow producers.

  \item \textbf{Channel Merge.} Write a function \lstinline|merge| that takes two channels and returns a new channel containing all values from both, interleaved in arrival order. Use \lstinline|select| in a loop with a \lstinline|default| arm to detect when both channels are drained.

  \item \textbf{Async Pipeline.} Using \lstinline|async_iter|, \lstinline|async_map|, and \lstinline|async_filter|, create a pipeline that generates the first 200 natural numbers, cubes each one, filters to keep only cubes less than 10,000, and prints the results.

  \item \textbf{Priority Select.} Implement a ``priority queue'' using two channels: a high-priority channel and a low-priority channel. Write a consumer loop that uses \lstinline|select| to prefer the high-priority channel but falls through to the low-priority one when the high-priority channel is empty. (Hint: use nested \lstinline|select| with \lstinline|default|.)

\end{enumerate}

%% ─────────────────────────────────────────────────────────────
\section*{What's Next}
%% ─────────────────────────────────────────────────────────────

You now have a solid grasp of channels---creation, send/recv semantics, closing, \lstinline|select| multiplexing, timeouts, and async iterators.
In the next chapter, we will step back and look at the bigger picture: when to use channels versus \lstinline|Ref| for shared state, how to avoid deadlocks, and a collection of battle-tested concurrency patterns including worker pools, broadcast, and request-reply.
